\documentclass[11pt]{article}
\usepackage[margin=1in]{geometry}
\usepackage{amsmath,amssymb,amsthm,mathtools}
\usepackage[hidelinks]{hyperref}
\allowdisplaybreaks[2]
\setlength{\emergencystretch}{3em}
\newtheorem{theorem}{Theorem}[section]
\newtheorem{lemma}[theorem]{Lemma}
\newtheorem{corollary}[theorem]{Corollary}
\theoremstyle{remark}\newtheorem{remark}[theorem]{Remark}
\newcommand{\E}{\mathbb E}
\newcommand{\dd}{\,d}
\newcommand{\one}{\mathbf1}
\newcommand{\Rhalf}{\mathcal R}
\title{Multiplicative strong-profile transfer in the overlap window}
\author{Analytical increment for review}
\date{Version 1; October 6, 2026}
\begin{document}\maketitle
\begin{abstract}
At the common right-half-circle arrival time, the original initialized
flow reproduces the exact target-only joint angular profile relative
to a same-field reference, with multiplicative error
$\exp[O(q^{10-2/\omega-1})]$. Radial weights have multiplicative
error $\exp[O(q^{10-2/\omega})]$. The proof normalizes the label
tangent by the target-only tangent. Its unperturbed generator is a
two-state contraction, preventing a logarithmic-time loss even for
arbitrarily small pair separations. The additive reference shift
is kept separate from shape. All results here concern only S1,
step 1; no strong-learning passage or LCRC entry at the later time
$t_0$ is claimed.
\end{abstract}

\section{Setting and the explicit overlap scales}
All new labels use prefix \texttt{an02owptv1:}.
Use the original balanced, untied initialized population of the
foundations and checkpoint P:
\[
 P_1=N(e_1,q^2I_2),\quad P_2=N(\rho e_2,q^2I_2),\quad
 P=(P_1+P_2)/2,\quad \rho\in[13/20,7/10],\quad S_0=q^{10}.
\]
The fixed initial label law is $\dd\alpha/(2\pi)$ and
$\theta(0,\alpha)=\psi(0,\alpha)=\alpha$, $m(0,\alpha)=S_0$.
Let $\Rhalf=[-\pi/2,\pi/2]$. Hats denote the exact aligned
target-only flow with the same initialization; they do not denote
the trained population.
Retain PROFILE's notation
\[
 \lambda=\rho^2,\quad \omega=(1-\lambda)/2,\quad
 a_\rho=\rho K(\rho a_\rho),\quad P_\rho=\Phi(\rho a_\rho),\quad
 d=(1-\lambda P_\rho)/2,\quad s=d/\omega,\quad p=3/s.
\]
Its exact target-only attracting angle is $qa_q$, with
$a_q=a_\rho+O(q^2)$. Fix a sufficiently large constant $C_m$
from PROFILE's uniform arrival/time guard, independent of all later
LCRC residual guards. Define
\begin{align}
 t_m&=\frac2\omega\log(1/q)+C_m,
 &\sigma&=10-\frac2\omega,
 &\zeta&=\sigma-1,\label{an02owptv1:time}\\
 \overline S(t)&=q^{10}e^{(1+2q^2)t},
 &\Lambda(t)&=\frac1q\int_0^t\overline S(v)\dd v.
 \label{an02owptv1:envelope}
\end{align}
The powers obey, by substituting the endpoints of the monotone
function $4/(1-\lambda)$,
\begin{equation}\label{an02owptv1:powers}
 \frac{110}{51}\le\sigma\le\frac{710}{231},\qquad
 \frac{59}{51}\le\zeta\le\frac{479}{231}.
\end{equation}
In particular, uniformly on the ratio box,
\begin{align}
 \overline S(t_m)&=e^{C_m}q^\sigma
                  \exp[2q^2t_m]=O(q^\sigma),
 \label{an02owptv1:massscale}\\
 \Lambda(t_m)&=
 \frac{\overline S(t_m)-q^{10}}{q(1+2q^2)}
 =O(q^\zeta)=o(1),\qquad q^2t_m=o(1).
 \label{an02owptv1:defectscale}
\end{align}
All constants below are uniform in $\rho$ for the fixed choice of
$C_m$. The source comparisons ensure $S(t)\le\overline S(t)$,
so for every fixed $\delta>0$, $S(t_m)/\delta\to0$.

\section{A multiplicative tangent estimate before and inside the chart}
Write $a_\theta=(\cos\theta,\sin\theta)$,
$b_\psi=(\cos\psi,\sin\psi)$, and use perpendicular unit vectors
with the usual positive orientation. At time $t$ fix the actual
population field and set
\[
 A_\theta=\E_P[XX^T\one_{\{a_\theta^TX>0\}}],\quad
 C_f(t,\theta)=\E_P[f_t(X)X^T\one_{\{a_\theta^TX>0\}}].
\]
Then $R=A_\theta-C_f$ and the angular receiver field is
\begin{equation}\label{an02owptv1:field}
 F(t,\theta,\psi)=
 \begin{pmatrix}b_\psi^TRa_\theta^\perp\\
 (b_\psi^\perp)^TRa_\theta\end{pmatrix}
 =F^0(\theta,\psi)-F^f(t,\theta,\psi).
\end{equation}
Here $F^0$ uses $A_\theta$ in place of $R$, and $F^f$ uses $C_f$.
The field is held fixed when differentiating a receiver or an initial
label: one is not varying the initial population measure.

\begin{lemma}[Receiver Jacobian comparison]
\label{an02owptv1:jacobian}
For $0\le t\le t_m$, uniformly in $\alpha\in\Rhalf$,
\begin{equation}\label{an02owptv1:jacobiandecomp}
 D F(t,\theta(t,\alpha),\psi(t,\alpha))
 =\begin{pmatrix}-g&h\\h&-g\end{pmatrix}+\mathcal E_\alpha(t),
 \qquad \|\mathcal E_\alpha(t)\|_\infty
 \le C\overline S(t)/q,
\end{equation}
where, at $\vartheta=\widehat\theta(t,\alpha)$,
\[
 g=a_\vartheta^TA_\vartheta a_\vartheta,\qquad
 h=(a_\vartheta^\perp)^TA_\vartheta a_\vartheta^\perp\ge0.
\]
The bound does not require that the receiver is already charted.
\end{lemma}
\begin{proof}
P's proved early comparison, applicable because
$\overline S(t_m)=o(1)$, gives
\begin{equation}\label{an02owptv1:early}
 |\theta-\widehat\theta|+|\psi-\widehat\theta|
 \le C\overline S(t),\qquad
 \left|\log\frac{m}{\widehat m}\right|\le4\overline S(t).
\end{equation}
PROFILE Proposition 4.1 applied to the whole population gives
$\|C_f\|\le a_*S$, $\|\partial_\theta C_f\|\le32S/q$,
where $a_*=1/2+q^2$. Its first, unscaled matrix bounds are global
in receiver angle. Differentiating \eqref{an02owptv1:field} therefore
gives $\|DF^f\|_\infty\le CS/q$ globally. In scaled angular
coordinates the velocity value is $O(S/q)$, and its receiver
derivative is still $O(S/q)$.

For the target matrix, $0\preceq A_\theta\preceq a_*I$ and
$A_\theta'a_\theta=0$. Also $\|A_\theta'\|\le32/q$:
the same two-endpoint polar formula used in PROFILE bounds each
degree-two angular weight by $16/q$, now with matrix integrand
$ee^T$. This proves the needed derivative bound without an
assumption on the input angle. Symmetry gives
$a_\theta^TA_\theta'=0$ as well.

Direct receiver differentiation gives
\[
 DF^0(\theta,\psi)=
 \begin{pmatrix}
 -b^TA_\theta a+b^TA_\theta'a^\perp & (b^\perp)^TA_\theta a^\perp\\
 (b^\perp)^TA_\theta a^\perp & -b^TA_\theta a
 \end{pmatrix}.
\]
The boundary term in the lower-left entry vanishes by $A_\theta'a=0$.
At $\theta=\psi=\vartheta$, the upper-left boundary term vanishes
too, and the matrix is precisely the one in
\eqref{an02owptv1:jacobiandecomp}. At the actual angles,
$b^TA_\theta'a^\perp=(b-a)^TA_\theta'a^\perp$ is bounded by
$C|\psi-\theta|/q$. All other differences are bounded by
$C(|\theta-\vartheta|+|\psi-\vartheta|)/q$, using the displayed
first derivative of $A$, not a second derivative. Combine these
bounds with \eqref{an02owptv1:early} and $S\le\overline S$.
The Gaussian boundary formula and continuity of the network on the
circle justify the receiver derivatives at fixed $q>0$.
\end{proof}

\begin{lemma}[Relative tangents and pairwise separations]
\label{an02owptv1:tangents}
There is a fixed $C$ such that, with $\eta(t)=C\Lambda(t)$ and
sufficiently small $q$, for every $0\le t\le t_m$,
\begin{equation}\label{an02owptv1:tangentrange}
 e^{-\eta(t)}\le
 \frac{\partial_\alpha\theta(t,\alpha)}{
       \partial_\alpha\widehat\theta(t,\alpha)},\quad
 \frac{\partial_\alpha\psi(t,\alpha)}{
       \partial_\alpha\widehat\theta(t,\alpha)}
 \le e^{\eta(t)}.
\end{equation}
Both ratios lie in the stated interval. For any $\alpha>\beta$
in $\Rhalf$, each of
\begin{equation}\label{an02owptv1:pairs}
 \frac{\theta(t,\alpha)-\theta(t,\beta)}{
       \widehat\theta(t,\alpha)-\widehat\theta(t,\beta)},\qquad
 \frac{\psi(t,\alpha)-\psi(t,\beta)}{
       \widehat\theta(t,\alpha)-\widehat\theta(t,\beta)}
\end{equation}
lies in $[e^{-\eta(t)},e^{\eta(t)}]$ as well.
\end{lemma}
\begin{proof}
Put $\widehat p=\partial_\alpha\widehat\theta$.
The exact target-only scalar field has
$v_q'=h-g$ (P's boundary identity), so
\[
 \widehat p(t,\alpha)=\exp\!\left(\int_0^t(h-g)\dd v\right)>0.
\]
For $p=(\partial_\alpha\theta,\partial_\alpha\psi)^T$,
differentiate the same-field characteristic with respect to its
initial label. The initial tangent is $(1,1)^T$, and
$p'=DF\,p$. Define $w=p/\widehat p$. By Lemma
\ref{an02owptv1:jacobian},
\begin{equation}\label{an02owptv1:normalized}
 w'=h\begin{pmatrix}-1&1\\1&-1\end{pmatrix}w+
 \mathcal E_\alpha w,\qquad w(0)=\boldsymbol1.
\end{equation}
The propagator of the first term has nonnegative entries, row sums
one, and operator norm one in $\ell^\infty$: explicitly its entries
are $(1\pm e^{-2\int_s^t h})/2$ in the two alternating positions.
It fixes $\boldsymbol1$. Variation of constants thus gives, with
$D(t)=\int_0^t\|\mathcal E_\alpha(v)\|_\infty\dd v$,
\[
 \|w(t)\|_\infty\le e^{D(t)},\qquad
 \|w(t)-\boldsymbol1\|_\infty\le e^{D(t)}-1.
\]
For $D(t)\le1/4$, the last bound is at most $2D(t)$ and
$1-2D(t)\ge e^{-4D(t)}$ (integrate
$2/(1-2u)\le4$ for $0\le u\le1/4$).
Hence every component lies between $e^{-4D(t)}$ and $e^{4D(t)}$.
Since $D(t)\le C\Lambda(t)=o(1)$ uniformly, this proves
\eqref{an02owptv1:tangentrange} after enlarging $C$.
Integrate the two positive derivative inequalities in the initial
label from $\beta$ to $\alpha$ to obtain
\eqref{an02owptv1:pairs}. There is no lower bound assumption on the
size of the pairwise separation. The denominator is positive by
scalar uniqueness and the displayed expression for $\widehat p$.

For fixed $q$, the receiver field is continuously differentiable
in its angular arguments and continuous in time on this compact
interval; the usual variational equation is therefore legitimate.
This differentiates labels in one fixed population, not the donor
field as the population varies. No sign of the original cross
derivatives is assumed: positivity is obtained from
\eqref{an02owptv1:normalized} and its small perturbation.
\end{proof}

\begin{remark}[Relation to the suggested $\int S/q$ budget]
\label{an02owptv1:actualmass}
The result has $\eta=O(\int\overline S/q)$ directly from the proved
early envelope. It also has $\eta=O(\int S/q)$ on this interval.
Indeed PROFILE's positive-coordinate lower growth gives
$\widehat U_1(t,\alpha)\ge\sqrt{S_0}\cos\alpha\,e^{t/2}$ on
$[0,\pi/4]$. Integrating its square and using
\eqref{an02owptv1:early},
\[
 S(t)\ge \frac1{16}e^{-4\overline S(t)}S_0e^t.
\]
Since $\sup_{t\le t_m}(4\overline S(t)+2q^2t)=o(1)$,
$\overline S(t)\le32 S(t)$ for small $q$. Thus the suggested
integrated scale is valid, without using an assumed future mass law.
The absolute comparison in \eqref{an02owptv1:early} is used to
estimate Jacobian coefficients, not as an additive shape error.
\end{remark}

\section{Overlap theorem and explicit transferred profiles}
Choose the reference label
\[
 \alpha_c=qa_q,\qquad
 (\theta_c(t),\psi_c(t))=(\theta(t,\alpha_c),\psi(t,\alpha_c)).
\]
This is an actual characteristic of the original field. Its target-only
comparison is stationary in angle: $\widehat\theta(t,\alpha_c)=qa_q$.
It is not forced to remain aligned in the original flow.
At time $t_m$, define
\[
 h_x(\alpha)=\frac{\theta(t_m,\alpha)-\theta_c(t_m)}q,
 \quad h_y(\alpha)=\frac{\psi(t_m,\alpha)-\psi_c(t_m)}q,
 \quad\widehat h(\alpha)=\widehat\theta(t_m,\alpha)/q-a_q,
 \quad D(\alpha)=|h_x(\alpha)|+|h_y(\alpha)|.
\]

\begin{theorem}[Exact original-flow overlap profile]
\label{an02owptv1:transfer}
Under only the initialized model and the proved PROFILE and early
comparison results used above, the following assertions hold for
all sufficiently small $q$.

Every label of $\Rhalf$ is in a bounded joint strong chart at $t_m$:
for a fixed $L$ from PROFILE,
\begin{equation}\label{an02owptv1:charted}
 |\theta(t_m,\alpha)|,|\psi(t_m,\alpha)|\le q(L+1),\qquad
 S(t_m)\le e^{C_m}q^\sigma e^{2q^2t_m}\ll\delta.
\end{equation}
The reference and shape errors satisfy
\begin{align}
 |\theta_c(t_m)/q-a_q|+|\psi_c(t_m)/q-a_q|
 &\le Cq^\zeta,\label{an02owptv1:reference}\\
 h_x(\alpha)&=\widehat h(\alpha)e^{\varepsilon_x(\alpha)},\quad
 h_y(\alpha)=\widehat h(\alpha)e^{\varepsilon_y(\alpha)},\quad
 |\varepsilon_x|,|\varepsilon_y|\le\eta_m:=Cq^\zeta,
 \label{an02owptv1:shape}\\
 m(t_m,\alpha)&=\widehat m(t_m,\alpha)e^{\varepsilon_m(\alpha)},
 \qquad |\varepsilon_m|\le r_m:=Cq^\sigma.
 \label{an02owptv1:weights}
\end{align}
At $\alpha_c$ both displacements are exactly zero; no ratio at zero
is needed. The pairwise version \eqref{an02owptv1:pairs} holds
for all other pairs, not only pairs involving the reference.

Put $c_0=\cos\alpha\ge0$. In PROFILE's two-sided Q1 region
$\alpha\in[\pi/4,\pi/2]$, both $h_x,h_y$ are positive and
\begin{equation}\label{an02owptv1:q1profile}
 h_x,h_y\asymp e^{-dC_m}q^s(c_0+q)^{-s}.
\end{equation}
On $[0,\pi/4]$ their absolute values are at most
$C e^{-dC_m}q^s$. On the Q4 region
$\alpha\in[-\pi/2,-\pi/4]$, both are negative and
\begin{equation}\label{an02owptv1:q4profile}
 -h_x,-h_y\asymp
 e^{-dC_m}q^{(1+\lambda)s}(c_0+q)^{-(1-\lambda)s}.
\end{equation}
On $[-\pi/4,0]$ the absolute values are at most
$C e^{-dC_m}q^{(1+\lambda)s}$.
On the entire right half-circle the radial profile is
\begin{equation}\label{an02owptv1:radialprofile}
 m(t_m,\alpha)\asymp e^{C_m}q^\sigma(c_0+q)^2,
 \qquad M_{\Rhalf}(t_m)\asymp e^{C_m}q^\sigma.
\end{equation}
The multiplicative error relative to each exact target-only quantity
is that in \eqref{an02owptv1:shape}--\eqref{an02owptv1:weights};
the power-law displays retain PROFILE's fixed comparison constants.
They do not assert that those constants become one.
\end{theorem}
\begin{proof}
PROFILE's entry theorem and right-half-circle corollary put all
target-only right-half-circle labels in $[-qL,qL]$ at $t_m$, after
fixing $C_m$ sufficiently large. Since
$\overline S(t_m)/q=O(q^\zeta)=o(1)$, the early comparison gives
\eqref{an02owptv1:charted} and \eqref{an02owptv1:reference}.
These are valid uses of an absolute error for a bounded chart and
the reference position, not for small relative shapes.

Apply the pairwise theorem to $\alpha$ and $\alpha_c$, reversing
their order when necessary. This proves the signed multiplicative
relations \eqref{an02owptv1:shape} even when the displacements are
much smaller than the reference shift. The radial estimate in
\eqref{an02owptv1:early} proves \eqref{an02owptv1:weights}.
The exact substitutions in PROFILE are
\begin{align*}
 \epsilon_m=q^{-s}e^{-dt_m}&=e^{-dC_m}q^s,\\
 s_-=2d=(1-\lambda)s,\qquad
 \epsilon_{-,m}=q^{-s_-}e^{-dt_m}
 &=e^{-dC_m}q^{(1+\lambda)s},\\
 S_0e^{t_m}&=e^{C_m}q^\sigma.
\end{align*}
Insert these into its already proved displacements and weights,
and use \eqref{an02owptv1:shape}--\eqref{an02owptv1:weights}.
Integrating the weight comparison over the fixed half-circle gives
the mass statement. No new entry or residual hypothesis has been
introduced.
\end{proof}

\begin{corollary}[Normalized radial profile and finite cutoff]
\label{an02owptv1:tail}
Let $\nu_m$ and $\widehat\nu_m$ be the original and target-only
radial probability measures on the initial labels $\Rhalf$, and put
$\widehat D=2|\widehat h|$. For every $z>0$,
\begin{equation}\label{an02owptv1:tailsandwich}
 e^{-2r_m}\widehat\nu_m\{\widehat D>e^{\eta_m}z\}
 \le\nu_m\{D>z\}
 \le e^{2r_m}\widehat\nu_m\{\widehat D>e^{-\eta_m}z\}.
\end{equation}
In particular, with $\epsilon_m=e^{-dC_m}q^s$ and $p=3/s$,
\begin{equation}\label{an02owptv1:tailpowers}
 \nu_m\{D>z\}\le C\min\{1,(\epsilon_m/z)^p\},
 \qquad
 \nu_m\{D>z\}\ge c(\epsilon_m/z)^p
 \quad(C_1\epsilon_m\le z\le c_1\epsilon_m q^{-s}).
\end{equation}
There is no tail beyond $C\epsilon_m q^{-s}=Ce^{-dC_m}$.
The Q4 cutoff is the smaller scale
$C\epsilon_{-,m}q^{-s_-}=Ce^{-dC_m}q^{2\lambda s}$,
and its own upper tail index is $3/[(1-\lambda)s]$.
\end{corollary}
\begin{proof}
Equation \eqref{an02owptv1:weights} also bounds the two normalizing
masses by $e^{\pm r_m}$, so the normalized density ratio on labels
lies in $[e^{-2r_m},e^{2r_m}]$.
Equation \eqref{an02owptv1:shape} gives
$e^{-\eta_m}\widehat D\le D\le e^{\eta_m}\widehat D$.
These two inequalities prove \eqref{an02owptv1:tailsandwich}.
Apply PROFILE's full-right-half-circle tail and cutoff result;
shrink and enlarge its fixed range constants once to accommodate
the factors tending to one. The Q4 claims follow from its separate
profile and $s_-=(1-\lambda)s$. This transfers a finite-cutoff
profile, not an untruncated asymptotic density.
\end{proof}

\section{Reference shift, scope, and stopping point}
\begin{remark}[Why the absolute angular comparison is insufficient]
\label{an02owptv1:absolutepitfall}
At $\rho=7/10$, $\zeta=59/51$. PROFILE's analytic bound
$P_\rho<153/250$ gives
\[
 s=\frac{1-(49/100)P_\rho}{51/100}
 >\frac{17503}{12750}
 >\frac{14750}{12750}=\frac{59}{51}.
\]
Thus the available scaled absolute-error budget $O(q^\zeta)$
decays more slowly than the core scale $q^s$. This does not claim a
lower bound on the actual angular error; it shows why that budget
cannot prove relative core accuracy. The tangent argument instead
places it in the coefficient defect and centers at the actual
same-field characteristic. The reference is $a_q+O(q^\zeta)$ in
each scaled coordinate, not identically $a_q$.
\end{remark}

This increment supplies the original strong family's joint profile,
radial weighting, full right-half-circle chart, and reference data
at $t_m$, with explicit powers. The mass there is still
$\asymp q^\sigma\ll1$. It does not identify $t_m$ with LCRC's later
passage entry $t_0$, prove strong learning or its guard-independent
integral $B_s$, or supply Q2 Cartesian entry, seed, transverse energy,
or secondary profile at $t_0$. Nor does it prove subsequent core or
chart retention, outside mass, or any S2 distance/gain contract.
The overlap chart constant is fixed by PROFILE and $C_m$, independent
of LCRC's $(B,M,K_Y)$. Its later continuation is still unproved.

The requested first link closes using the already proved sources
and the displayed variational argument. Work stops here: no
$M$-clock passage, capped radial-gain proof, bridge argument, or
consolidation is performed.

\begin{thebibliography}{9}\small
\bibitem{PROFILE} PROFILE,
\emph{AN02 strong entry profile and tail budget v1}.
Prefix \texttt{inc:AN02:profile:v1:}; labels \texttt{root},
\texttt{arrival}, \texttt{entry}, \texttt{rightcircle}, \texttt{cost}.
\bibitem{P} P, \emph{THEOREM A ANALYTICAL PROGRESS}, checkpoint 1.0.
Labels \texttt{pa:earlybounds}, \texttt{pa:targetcomparison},
\texttt{pa:finiteqorder}; only the exact analytic portions are used.
\bibitem{F} \emph{RELU population foundations v2}, exact balanced
characteristic equations, Gaussian gate-boundary regularity and P5.
\bibitem{LCRC} LCRC, \emph{AN03 ledger cap residual closure v1}.
Section 1 and Theorem 4.1 define the later contracts; none is used
as a premise to this overlap proof.
\bibitem{ROADMAP} \emph{Unified analytic theory roadmap}, v1.2,
October 6, 2026: S0 retained inputs and assembly requirements;
S1, step 1.
\end{thebibliography}
\end{document}
